\documentclass[varwidth=16cm, border=2mm]{standalone}
\usepackage{amsmath}\usepackage{amssymb}\usepackage{ctex}
\usepackage{tasks}\usepackage{xcolor}\usepackage{graphicx}
\settasks{label=\Alph*., label-width=1.5em, item-indent=2em}
\begin{document}
\textbf{[0.6]} (2024·全国·高三对口高考) 已知椭圆 $\frac{x^2}{9} + y^2 = 1$，过左焦点 $F$ 作倾斜角为 $\frac{\pi}{6}$ 的直线交椭圆于 $A, B$ 两点，则弦 $AB$ 的长为\underline{\hspace{1.5cm}}.

\vspace{0.2cm}\par\noindent\textcolor{red}{\textbf{【答案】} 2}
\par\noindent\textcolor{blue}{\textbf{【解析】} 1. **识别椭圆方程参数**：\n   已知椭圆方程为 $\frac{x^2}{9} + y^2 = 1$。\n   根据标准方程 $\frac{x^2}{a^2} + \frac{y^2}{b^2} = 1$，可得 $a^2 = 9 \Rightarrow a = 3$， $b^2 = 1 \Rightarrow b = 1$。\n   椭圆的半焦距 $c$ 满足 $c^2 = a^2 - b^2$，所以 $c^2 = 9 - 1 = 8 \Rightarrow c = \sqrt{8} = 2\sqrt{2}$。\n   左焦点 $F$ 的坐标为 $(-c, 0) = (-2\sqrt{2}, 0)$。\n\n2. **确定直线方程**：\n   直线过左焦点 $F(-2\sqrt{2}, 0)$，倾斜角为 $\frac{\pi}{6}$。\n   直线的斜率 $k = \tan(\frac{\pi}{6}) = \frac{\sqrt{3}}{3}$。\n   根据点斜式，直线方程为 $y - 0 = \frac{\sqrt{3}}{3}(x - (-2\sqrt{2}))$，即 $y = \frac{\sqrt{3}}{3}(x + 2\sqrt{2})$。\n\n3. **计算弦长 (方法一：韦达定理与弦长公式)**：\n   将直线方程代入椭圆方程：\n   $\frac{x^2}{9} + \left(\frac{\sqrt{3}}{3}(x + 2\sqrt{2})\right)^2 = 1$ \n   $\frac{x^2}{9} + \frac{3}{9}(x + 2\sqrt{2})^2 = 1$ \n   $\frac{x^2}{9} + \frac{1}{3}(x^2 + 4\sqrt{2}x + 8) = 1$ \n   两边同乘以 $9$ 整理：\n   $x^2 + 3(x^2 + 4\sqrt{2}x + 8) = 9$ \n   $x^2 + 3x^2 + 12\sqrt{2}x + 24 = 9$ \n   $4x^2 + 12\sqrt{2}x + 15 = 0$ \n   设交点 $A, B$ 的横坐标分别为 $x_1, x_2$。根据韦达定理：\n   $x_1 + x_2 = -\frac{12\sqrt{2}}{4} = -3\sqrt{2}$ \n   $x_1 x_2 = \frac{15}{4}$ \n   弦长公式为 $|AB| = \sqrt{(x_1 - x_2)^2 + (y_1 - y_2)^2}$。\n   由于 $y = k(x + 2\sqrt{2})$，所以 $y_1 - y_2 = k(x_1 - x_2)$。\n   因此，弦长 $|AB| = \sqrt{(x_1 - x_2)^2 + k^2(x_1 - x_2)^2} = \sqrt{(x_1 - x_2)^2(1 + k^2)} = |x_1 - x_2|\sqrt{1 + k^2}$。\n   首先计算 $(x_1 - x_2)^2 = (x_1 + x_2)^2 - 4x_1 x_2 = (-3\sqrt{2})^2 - 4(\frac{15}{4}) = (9 \times 2) - 15 = 18 - 15 = 3$。\n   所以 $|x_1 - x_2| = \sqrt{3}$。\n   又因为 $k = \frac{\sqrt{3}}{3}$，所以 $k^2 = (\frac{\sqrt{3}}{3})^2 = \frac{3}{9} = \frac{1}{3}$。\n   $1 + k^2 = 1 + \frac{1}{3} = \frac{4}{3}$。\n   将这些值代入弦长公式：\n   $|AB| = \sqrt{3} \times \sqrt{\frac{4}{3}} = \sqrt{3} \times \frac{2}{\sqrt{3}} = 2$。\n\n4. **计算弦长 (方法二：椭圆焦弦公式)**：\n   椭圆的离心率 $e = \frac{c}{a} = \frac{2\sqrt{2}}{3}$。\n   过焦点的弦长公式为 $|AB| = \frac{2a(1 - e^2)}{1 - e^2 \cos^2 \theta}$，其中 $\theta$ 为焦弦与长轴的夹角，即直线的倾斜角 $\frac{\pi}{6}$。\n   $1 - e^2 = 1 - (\frac{2\sqrt{2}}{3})^2 = 1 - \frac{8}{9} = \frac{1}{9}$。\n   $\cos \theta = \cos(\frac{\pi}{6}) = \frac{\sqrt{3}}{2}$，所以 $\cos^2 \theta = (\frac{\sqrt{3}}{2})^2 = \frac{3}{4}$。\n   将这些值代入焦弦公式：\n   $|AB| = \frac{2 \times 3 \times (\frac{1}{9})}{1 - \frac{8}{9} \times \frac{3}{4}} = \frac{\frac{6}{9}}{1 - \frac{2}{3}} = \frac{\frac{2}{3}}{\frac{1}{3}} = 2$。\n\n两种方法计算结果均为 $2$。}


\end{document}
